%% 自動生成（tools/mkchecklist.py）。各章の「ふりかえり」の項目を集めたもの
\begingroup\footnotesize
\def\bkcb{\tikz[baseline=-0.5ex]\draw[\bkTF, line width=0.7pt, rounded corners=0.5pt] (0,-0.3em) rectangle (0.8em,0.5em);}
\begin{longtblr}{colspec={Q[c,m,wd=9mm] X[l,m] Q[c,m,wd=10mm] Q[c,m,wd=9mm] Q[c,m,wd=9mm]}, width=\linewidth, rowhead=1, colsep=3pt,
  row{1}={bg=bkVol!14, font=\sffamily\bfseries\scriptsize}, hlines={0.4pt,bkLine}, rowsep=0.7pt}
章 & できるようになったこと & ページ & 1回め & 2回め \\
{\sffamily\bfseries\color{\bkTT}8-1} & 地図の等高線や天気図の等圧線から、高い所・低い所・急な所を読みとれる。 & \pageref{bk:8-1:kankaku} & \bkcb & \bkcb \\
 & 2つの変数の関数を、表・鳥かん図・等高線の3通りで表し、表計算と Python で作れる。 & \pageref{bk:8-1:hyou} & \bkcb & \bkcb \\
 & 等高線の式と切り口の関数を求め、3つの変数の関数の等値面を説明できる。 & \pageref{bk:8-1:daen} & \bkcb & \bkcb \\
{\sffamily\bfseries\color{\bkTT}8-2} & 等高線や数の表から、東の坂と北の坂を別々に読みとれる。 & \pageref{bk:8-2:hyou} & \bkcb & \bkcb \\
 & ほかを止めて1つだけ動かす偏微分を計算し、点での値の意味を言える。 & \pageref{bk:8-2:henbibun} & \bkcb & \bkcb \\
 & 2回の偏微分を計算して順番によらないことを確かめ、3つの変数の量に使える。 & \pageref{bk:8-2:nikai} & \bkcb & \bkcb \\
{\sffamily\bfseries\color{\bkTT}8-3} & ななめの変化を、東の変化と北の変化のたし算で見積もれる。 & \pageref{bk:8-3:kaidan} & \bkcb & \bkcb \\
 & 全微分で、近似の値と誤差の伝わり方を計算できる。 & \pageref{bk:8-3:zenbibun} & \bkcb & \bkcb \\
 & 連鎖律とヤコビ行列で、手の動きと面積の倍率（ヤコビアン）を求められる。 & \pageref{bk:8-3:jacobi} & \bkcb & \bkcb \\
{\sffamily\bfseries\color{\bkTT}8-4} & 等高線の地図に、ボールが転がりはじめる向き（等高線と直角）をかきこめる。 & \pageref{bk:8-4:chokkaku} & \bkcb & \bkcb \\
 & 東の坂と北の坂を並べて勾配 grad を作り、どの向きの坂も内積で計算できる。 & \pageref{bk:8-4:houkou} & \bkcb & \bkcb \\
 & 内積の $\cos$ の形から、いちばん急な向きと等高線の向きを説明し、$\vec F=-\mathrm{grad}\,U$ などの物理の式を使える。 & \pageref{bk:8-4:cos} & \bkcb & \bkcb \\
{\sffamily\bfseries\color{\bkTT}8-5} & 場所ごとの矢印を計算して、縮める割合をそろえた矢印の場の図にかける。 & \pageref{bk:8-5:chijimi} & \bkcb & \bkcb \\
 & 主役の6つの場・点電荷の場・勾配の場を見分け、流線（力線）の式を求められる。 & \pageref{bk:8-5:ryusenshiki} & \bkcb & \bkcb \\
 & 座標の軸を回して見たときの矢印の成分を、回転行列の転置で計算できる。 & \pageref{bk:8-5:kaiten} & \bkcb & \bkcb \\
{\sffamily\bfseries\color{\bkTT}8-6} & 板を小さなタイルに分けて重さを表計算でたし、重積分の記号を「何を・どこで・1枚の量」で読める。 & \pageref{bk:8-6:tile} & \bkcb & \bkcb \\
 & 長方形の範囲の重積分を、ふつうの定積分を2回くり返す形に直して計算できる。 & \pageref{bk:8-6:naosu} & \bkcb & \bkcb \\
 & 場合分けの表を使って、三角形・円・立体の重積分を、正しいふつうの定積分に直して計算できる。 & \pageref{bk:8-6:basho} & \bkcb & \bkcb \\
{\sffamily\bfseries\color{\bkTT}8-7} & 道を小さな1歩に分け、力と1歩の内積をたして仕事を表計算で求め、線積分の記号を読める。 & \pageref{bk:8-7:ippo} & \bkcb & \bkcb \\
 & 軸に平行な道（折れ線）の線積分を、ふつうの定積分に直して計算できる。 & \pageref{bk:8-7:jiku} & \bkcb & \bkcb \\
 & 場合分けの表を使って、ななめ・曲線・円の道、保存力、閉じた道の線積分を計算できる。 & \pageref{bk:8-7:basho} & \bkcb & \bkcb \\
{\sffamily\bfseries\color{\bkTT}8-8} & 網を通る水の量が、流れの「面に直角な成分」×面積で決まることを説明できる。 & \pageref{bk:8-8:chokkaku} & \bkcb & \bkcb \\
 & 平らな面の流束を、法線ベクトルの内積と重積分で求められる。 & \pageref{bk:8-8:taira} & \bkcb & \bkcb \\
 & 場合分けの表で、曲がった面や閉じた面の面積分をふつうの定積分に直せる。 & \pageref{bk:8-8:basho} & \bkcb & \bkcb \\
{\sffamily\bfseries\color{\bkTT}8-9} & 矢印の図にマスを置いて、わき出しと吸いこみを見分けられる。 & \pageref{bk:8-9:hako} & \bkcb & \bkcb \\
 & マスの辺を通る量から、表計算で発散を求められる。 & \pageref{bk:8-9:masu} & \bkcb & \bkcb \\
 & 発散を偏微分の和で計算し、流束や連続の式と結びつけられる。 & \pageref{bk:8-9:shiki} & \bkcb & \bkcb \\
{\sffamily\bfseries\color{\bkTT}8-10} & 矢印の場の図から、小さな風車が回るかどうかと、その向きを見分けられる。 & \pageref{bk:8-10:fusha} & \bkcb & \bkcb \\
 & マスのまわりを1周してたす量（循環）を表計算で求め、面積でわって回転を見積もれる。 & \pageref{bk:8-10:masu} & \bkcb & \bkcb \\
 & 平面と3次元の rot を偏微分で計算し、保存力の見分けや、電流のまわりの磁場に使える。 & \pageref{bk:8-10:heimen} & \bkcb & \bkcb \\
{\sffamily\bfseries\color{\bkTT}8-11} & マスごとのわき出し（うず）を全部たすと、外のふちから出る量（外の1周）になることを、表計算で確かめられる。 & \pageref{bk:8-11:uchikeshi} & \bkcb & \bkcb \\
 & 閉じた面や閉じた道の積分を、2つの定理で「計算が楽な側」の積分に取りかえて求められる。 & \pageref{bk:8-11:gauss} & \bkcb & \bkcb \\
 & 対称性に合った袋や道を選び、ガウスの法則・アンペールの法則で、場の大きさを求められる。 & \pageref{bk:8-11:hosoku} & \bkcb & \bkcb \\
{\sffamily\bfseries\color{\bkTT}8-12} & マスの値と「まわりの平均」との差を出し、温まる所と冷える所を見分けられる。 & \pageref{bk:8-12:sa} & \bkcb & \bkcb \\
 & 「どのマスもまわりの平均と同じ」のくり返しで、板の落ち着いた温度を求められる。 & \pageref{bk:8-12:ita} & \bkcb & \bkcb \\
 & ラプラシアンを偏微分で計算し、球や円柱の形の場合にも使える。 & \pageref{bk:8-12:lap} & \bkcb & \bkcb \\
{\sffamily\bfseries\color{\bkTT}8-13} & 棒の温度を「となりとの差」の表計算で計算できる。 & \pageref{bk:8-13:masu} & \bkcb & \bkcb \\
 & こわれない条件 $r\leqq\frac12$ を1歩の倍率で説明できる。 & \pageref{bk:8-13:joken} & \bkcb & \bkcb \\
 & 熱と波の方程式を作り、代入して確かめられる。 & \pageref{bk:8-13:netsu} & \bkcb & \bkcb \\
{\sffamily\bfseries\color{\bkTT}8-14} & 4つの式を、「わき出し」と「うず」のことばで読める。 & \pageref{bk:8-14:yomu} & \bkcb & \bkcb \\
 & 積分形と微分形のつながり、光の速さの出し方を言える。 & \pageref{bk:8-14:sekibun} & \bkcb & \bkcb \\
 & 流れの式・熱の式の項を、この巻のことばに訳せる。 & \pageref{bk:8-14:ns} & \bkcb & \bkcb \\
\end{longtblr}
\endgroup
